\documentclass[12pt]{letter}
\usepackage[margin=1in]{geometry}
\usepackage{hyperref}

\signature{Anonymous Authors}
\address{Department of Computer Science\\Anonymous University}

\begin{document}

\begin{letter}{Editorial Office\\Information Fusion\\Elsevier}

\opening{Dear Editor,}

We are pleased to submit our manuscript entitled ``\textbf{Multi-Source Retrieval-Augmented Generation for Personal Knowledge Management: Integrating Co-occurrence Statistics, Knowledge Graphs, and Ontological Reasoning via Learned Gating}'' for consideration in \textit{Information Fusion}.

\textbf{Relevance to Information Fusion.} Our work addresses a core information fusion challenge: how to optimally combine heterogeneous information sources---dense vector retrieval, statistical co-occurrence, and knowledge graph reasoning---for retrieval-augmented generation. The learned gating mechanism we propose is directly analogous to multi-sensor fusion, dynamically weighting source contributions based on query characteristics. This falls squarely within the scope of Information Fusion's interest in novel fusion architectures and adaptive combination strategies.

\textbf{Key contributions:}
\begin{enumerate}
    \item \textbf{Multi-source fusion architecture:} We are the first to systematically integrate all three retrieval paradigms (dense, statistical, graph-based) with OWL ontological reasoning for personal knowledge management.
    
    \item \textbf{Learned gating mechanism:} Our sigmoid-based gating network achieves 18.4\% improvement over heuristic reciprocal rank fusion on HotpotQA ($p<0.001$, Cohen's $d=0.475$), demonstrating the value of query-adaptive fusion.
    
    \item \textbf{Dual-benchmark validation:} Evaluation on both HotpotQA and MuSiQue shows consistent multi-source advantages, with larger gains on harder compositional queries---confirming that source fusion is most valuable when individual sources are insufficient.
    
    \item \textbf{Statistical rigor:} All results include confidence intervals, effect sizes, and significance tests.
\end{enumerate}

\textbf{Novelty.} While prior work has explored pairwise combinations (dense + sparse retrieval, or dense + knowledge graph), no existing framework combines all three paradigms with ontological reasoning and learned fusion. Our per-query-type analysis provides actionable insights into when each source dominates, contributing to the theoretical understanding of information fusion for retrieval.

The manuscript has not been published or submitted elsewhere. All authors have approved the submission.

\closing{Sincerely,}

\end{letter}
\end{document}
